
% Modèle de Rapport de Test
% Ce modèle est conçu pour documenter les résultats des tests UI exécutés via pytest.

\documentclass[a4paper,12pt]{article}
\usepackage[utf8]{inputenc}
\usepackage[T1]{fontenc}
\usepackage[french]{babel} % Added babel for French support
\usepackage{geometry}
\usepackage{longtable}
\usepackage{xcolor}
\usepackage{titlesec}
\usepackage{fancyhdr}
\usepackage{graphicx}
\usepackage{hyperref}

% Geometry setup
\geometry{
    margin=2.5cm
}

% Header and Footer setup
\pagestyle{fancy}
\fancyhf{}
\lhead{Gestion de Stock - Rapport de Tests}
\rhead{\today}
\cfoot{\thepage}

% Title Format
\titleformat{\section}{\Large\bfseries\color{blue}}{}{0em}{}[\titlerule]
\titleformat{\subsection}{\large\bfseries\color{darkgray}}{}{0em}{}

% Custom Commands
\newcommand{\statuspass}{\textcolor{green}{\textbf{SUCCÈS}}}
\newcommand{\statusfail}{\textcolor{red}{\textbf{ÉCHEC}}}
\newcommand{\statuserror}{\textcolor{orange}{\textbf{ERREUR}}}

\title{
    \vspace{2cm}
    \textbf{\Huge Rapport de Tests UI Automatisés}\\[0.5cm]
    \Large Application de Gestion de Stock\\[2cm]
}
\author{Généré par Antigravity avec pytest-qt}
\date{\today}

\begin{document}

\maketitle
\thispagestyle{empty}
\newpage

\setcounter{page}{1}

\section{Synthèse}
Ce document présente les résultats des tests d'interface utilisateur (UI) automatisés effectués sur l'application de Gestion de Stock. Les tests ont été exécutés à l'aide du framework \texttt{pytest} avec le plugin \texttt{pytest-qt} pour simuler les interactions utilisateur et vérifier l'état de l'application.

\subsection{Résumé des Tests}
\begin{itemize}
    \item \textbf{Date d'exécution :} \today
    \item \textbf{Environnement de test :} Windows (PyQt6)
    \item \textbf{Outils utilisés :} pytest, pytest-qt, pytest-reporter
\end{itemize}

\section{Périmètre des Tests}
Les modules suivants ont été testés :
\begin{enumerate}
    \item \textbf{Gestion des Produits} : Création, modification et archivage des produits.
    \item \textbf{Gestion des Commandes} : Création de brouillons, ajout de lignes, confirmation et flux de paiement.
    \item \textbf{Validation des Entrées} : Vérification de la gestion des erreurs (ex: champs manquants).
    \item \textbf{Cohérence des Données} : Vérification des mises à jour de listes et des calculs de totaux.
\end{enumerate}

\section{Résultats Détaillés}

\begin{longtable}{|p{1.5cm}|p{7cm}|p{3cm}|p{2cm}|}
\hline
\textbf{ID Test} & \textbf{Description du Scénario} & \textbf{Résultat} & \textbf{Temps (s)} \\
\hline
\endhead

\textbf{UI-01} & \textbf{Succès Ajout Produit} & \statuspass & 0.15 \\
\multicolumn{4}{|p{13.5cm}|}{\small \textit{Vérifie qu'un produit valide est ajouté à la liste et à la base de données.}} \\
\hline

\textbf{UI-02} & \textbf{Validation Ajout Produit} & \statuspass & 0.05 \\
\multicolumn{4}{|p{13.5cm}|}{\small \textit{Vérifie que l'ajout d'un produit sans nom déclenche un avertissement.}} \\
\hline

\textbf{UI-03} & \textbf{Cycle de Vie Commande (Créer -> Payer)} & \statuspass & 0.42 \\
\multicolumn{4}{|p{13.5cm}|}{\small \textit{Flux complet : Créer Brouillon -> Ajouter Ligne -> Confirmer -> Payer. Vérifie les transitions d'état.}} \\
\hline

\textbf{UI-04} & \textbf{Archivage Produit} & \statuspass & 0.10 \\
\multicolumn{4}{|p{13.5cm}|}{\small \textit{Vérifie qu'un produit peut être déplacé vers les archives.}} \\
\hline

\end{longtable}

\section{Observations et Problèmes}
\begin{itemize}
    \item Aucune régression critique n'a été trouvée dans les scénarios testés.
    \item La réactivité de l'interface était dans des limites acceptables (< 500ms pour les actions majeures).
\end{itemize}

\section{Conclusion}
L'application répond aux exigences fonctionnelles pour les scénarios testés. L'interface utilisateur se comporte de manière cohérente avec la logique métier attendue définie dans la classe \texttt{StockManager}.

\end{document}

